Une charge $q_1 = +\qty{\v{q1}}{\micro C}$ et une charge $q_2 = -\qty{\v{q2}}{\micro C}$ sont placées à \qty{\v{r}}{cm} l'une de l'autre. La constante de Coulomb est $k = \qty{8.99e9}{N.m^2/C^2}$.

\begin{enumerate}[label=\alph*)]
    \item Déterminez la grandeur de la force électrique entre les deux charges.
    \item Cette force est-elle attractive ou répulsive ? Justifiez.
    \item À quelle distance faudrait-il placer les charges pour que la force soit de \qty{\v{Fc}}{N} ?
\end{enumerate}
